\section{Round-five positive closure: the weighted gauge complex, Radon duality, and the joint Gaussian process}
\label{sec:r5-b3}

The collision increment and collision source are placed in the same weighted
space.  The exponential Hamiltonian is used only on an open integrability
chart.  The covariance kernel is identified by local hard-sphere variations,
and process convergence is proved with Aldous--Mitoma estimates including the
finite-particle jump term.

\subsection{Ambient weighted spaces and the exponential chart}

Put
\[
 w_p(v)=1+|v|^2,
 \qquad w_c(v,v_*)=1+|v|^2+|v_*|^2.
\]
Let \(\mathcal Y_p\) be the completion of smooth tests for the norm
\[
 \|p\|_{\mathcal Y_p}=
 \left\|\frac{p_0}{w_p}\right\|_\infty+
 \left\|\frac{p_T}{w_p}\right\|_\infty+
 \left\|\frac{(\partial_t+v\cdot\nabla_x)p}{w_p}\right\|_\infty+
 \left\|\frac{\Delta_\omega p}{w_c}\right\|_\infty.
\]
Let \(\mathcal Y_\psi=C_{w_c}\) be the weighted continuous collision-source
space and put \(\mathcal Y_z=C_{w_c}\).  The ambient linear map is
\[
 \mathfrak D(p,\psi)=\Delta p+\psi.
\]
For a Maxwellian exponent \(\beta>0\), choose \(0<\eta<\beta/8\) and define
the open Hamiltonian chart
\[
 \mathcal U_\eta=\left\{(p,\psi):
 \|\mathfrak D(p,\psi)/w_c\|_\infty<\eta\right\}.
\]
On energy-moment sublevels the integral of
\(e^{\mathfrak D(p,\psi)}\) is finite on this chart.

\begin{lemma}[Typing of the collision increment]
\label{lem:r5-b3-typing}
The maps
\[
 \Delta:\mathcal Y_p\to\mathcal Y_z,
 \qquad
 \mathfrak D:\mathcal Y_p\times\mathcal Y_\psi\to\mathcal Y_z
\]
are continuous.  The Hamiltonian
\[
 \mathbb H(f,p,\psi)=
 \langle f,v\cdot\nabla_xp\rangle+
 \int(e^{\Delta p+\psi}-1)dA_f
\]
is finite and analytic on \(\mathcal U_\eta\), locally uniformly on the
Maxwellian moment sublevels supplied by B2.
\end{lemma}

\begin{proof}
Continuity of \(\Delta\) is part of the declared norm, so no quadratic-growth
increment is placed in an unweighted bounded space.  If
\(|z|\le\eta w_c\), then the collision flux times \(e^z\) is dominated by a
product of Maxwellians with exponent \(\beta-4\eta>0\), together with a
polynomial relative-velocity factor.  This is integrable uniformly on the
moment sublevels.  The same majorant multiplied by any power of \(z\) proves
analyticity by dominated convergence.
\end{proof}

\subsection{Closed balance and the exact gauge}

For a pair \((f,\Gamma)\), define its balance defect
\[
 \mathfrak B_{f,\Gamma}(p)=
 \langle p_T,f_T\rangle-\langle p_0,f_0\rangle
 -\int_0^T\langle(\partial_t+v\cdot\nabla_x)p,f\rangle dt
 -\int\Delta p\,d\Gamma.
\]
The pair is balanced when this vanishes for every \(p\in\mathcal Y_p\).

\begin{lemma}[Closed weighted balance]
\label{lem:r5-b3-balance}
On sets with bounded second moment and contact entropy, the relation
\(\mathfrak B_{f,\Gamma}=0\) is closed in the weighted weak topology.
\end{lemma}

\begin{proof}
Truncate velocities at \(R\).  Every truncated term is continuous.  Energy
bounds the free-transport tail.  For the collision tail, Young's inequality
\(ab\le e^a+\ell(b)\), with
\(a=cw_c\), gives uniform integrability of \(w_c\Gamma\) from the Maxwellian
reference and the entropy bound.  Let \(R\to\infty\) after taking the weak
limit.
\end{proof}

Define
\[
 \mathfrak G r=(r,-\Delta r).
\]

\begin{theorem}[Exact weighted gauge complex]
\label{thm:r5-b3-gauge}
The sequence
\[
 0\longrightarrow\mathcal Y_p
 \xrightarrow{\mathfrak G}
 \mathcal Y_p\times\mathcal Y_\psi
 \xrightarrow{\mathfrak D}
 \mathcal Y_z\longrightarrow0
\]
is exact.  The range of \(\mathfrak G\) is closed and
\[
 (\mathcal Y_p\times\mathcal Y_\psi)/\operatorname{Ran}\mathfrak G
 \cong\mathcal Y_z
\]
as Hausdorff Banach spaces.  The open Hamiltonian chart is the inverse image
under this quotient map of the open ball
\(\{z:\|z/w_c\|_\infty<\eta\}\).
\end{theorem}

\begin{proof}
Surjectivity follows from \(z=\mathfrak D(0,z)\).  If
\(\mathfrak D(p,\psi)=0\), then
\((p,\psi)=(p,-\Delta p)=\mathfrak Gp\).  The range is the graph of the
continuous map \(-\Delta\), hence is closed.  The induced quotient map has
continuous inverse \(z\mapsto[(0,z)]\).  The final chart statement is then
immediate.
\end{proof}

\subsection{Weighted-strict Radon duality}

Let \(C_{w_c}\) carry the weighted bounded-strict topology \(\beta_{w_c}\):
a net converges when it is bounded in \(\|z/w_c\|_\infty\) and converges
uniformly on every compact weight sublevel.  Let
\(\mathcal M_{w_c}\) be the countably additive signed Radon measures with
finite \(w_c\)-moment.

\begin{lemma}[No finitely additive collision dual]
\label{lem:r5-b3-radon}
The continuous dual of \((C_{w_c},\beta_{w_c})\) is exactly
\(\mathcal M_{w_c}\) under integration.
\end{lemma}

\begin{proof}
A continuous functional restricted to a compact weight sublevel is sup-norm
continuous and hence has a Radon representation.  Compatibility of the
representations follows from restriction.  Bounded-set continuity makes the
weighted total variations uniformly summable, so the compact pieces paste to
one countably additive measure with finite \(w_c\)-moment.  Conversely such a
measure is continuous by tightness of \(w_c|m|\) and dominated convergence.
\end{proof}

For fixed \(f\), put
\[
 \mathfrak J_f(\Gamma)=
 \begin{cases}
 \displaystyle\int\ell(d\Gamma/dA_f)dA_f,&\Gamma\ll A_f,\\
 +\infty,&\text{otherwise}.
 \end{cases}
\]

\begin{theorem}[Closed path-space Fenchel duality]
\label{thm:r5-b3-duality}
For every finite-moment pair,
\[
 \sup_{(p,\psi)\in\mathcal U_\eta\text{ for some }\eta<\beta/8}
 \left\{
 \begin{aligned}
 &\langle p_T,f_T\rangle-\langle p_0,f_0\rangle
 -\int\langle(\partial_t+v\cdot\nabla)p,f\rangle dt\\
 &\quad+\int\psi\,d\Gamma
 -\int(e^{\Delta p+\psi}-1)dA_f
 \end{aligned}
 \right\}
 =
 \begin{cases}
 \mathfrak J_f(\Gamma),&(f,\Gamma)\text{ balanced},\\
 +\infty,&\text{otherwise}.
 \end{cases}
\]
The supremum over the increasing integrability charts is unchanged by
weighted-strict closure of smooth compactly supported sources.
\end{theorem}

\begin{proof}
If balance fails, Lemma~\ref{lem:r5-b3-balance} and Hahn--Banach give a test
\(p\) with nonzero defect.  Take \(\psi=-\Delta p\) and scale first within a
small chart; repeating the argument on increasing charts makes the supremum
infinite.  On a balanced pair the expression depends only on
\(z=\Delta p+\psi\) and becomes
\[
 \int z\,d\Gamma-\int(e^z-1)dA_f.
\]
Because \(p=0\) is allowed, every weighted source in every chart occurs.
Scalar conjugacy gives \(\ell\) for bounded density ratios.  Monotone
truncation of \(\log(d\Gamma/dA_f)\) handles an unbounded absolutely continuous
part.  For a singular part use increasing compact sets and sources
\(z=n1_K\), smoothed in \(\beta_{w_c}\).  Lemma~\ref{lem:r5-b3-radon}
justifies all pairings with countably additive measures.
\end{proof}

\subsection{The complete covariance kernel}

A full linear source is a quadruple of initial, terminal, bulk and collision
parts.  The balance adjoint maps \(p\) to
\[
 \mathfrak A^*p=(-p_0,p_T,-(\partial_t+v\cdot\nabla)p,-\Delta p).
\]
Collision invariants are included as the finite-dimensional kernel of
\(\Delta\).

\begin{theorem}[Hard-sphere cohomological kernel]
\label{thm:r5-b3-kernel}
On a regular B2 source chart, a real source direction has zero limiting
speed-normalized variance if and only if it belongs to the closed span of
\(\operatorname{Ran}\mathfrak A^*\) and the deterministic total mass,
momentum and energy directions.  Equivalently the pressure Hessian is strictly
positive on every finite-dimensional family of independent quotient classes.
\end{theorem}

\begin{proof}
Every adjoint direction is the exact microscopic balance identity and every
collision invariant is deterministic under the corresponding prepared shell,
so these directions have zero quotient variance.

Conversely suppose a source has zero variance.  B2 normal convergence makes
its microscopic variance vanish on every open regular phase configuration.
First vary one initial point along a collision-free free-flight cylinder.  The
fundamental lemma of the calculus of variations gives a potential \(p\) whose
transport derivative equals the bulk source and whose traces equal the
endpoint sources, modulo a function of the conserved totals.  Next vary the
impact normal and the two incoming velocities in an open set containing
exactly one binary contact.  The contact-source difference must equal
\(-\Delta p\); otherwise its value changes on that open contact cylinder and
has positive variance.  Subtracting \(\mathfrak A^*p\) leaves a function
constant on all collision-free and one-collision open sets.  Connectivity by
successive binary collisions implies that this remainder is a linear
combination of total mass, momentum and energy.  Approximation in the weighted
source norm gives the closed-range statement.
\end{proof}

\subsection{The joint Gaussian martingale problem}

Let \((f,\Gamma=qA_f)\) be a regular biased law and define
\[
 \zeta^\varepsilon=\sqrt{\mu_\varepsilon}(\pi^\varepsilon-f),
 \qquad
 \Xi^\varepsilon=\sqrt{\mu_\varepsilon}
 (\Gamma^\varepsilon-\Gamma).
\]
For density test \(\phi\) and contact test \(\psi\), let
\(M^\varepsilon(\phi)\) and \(N^\varepsilon(\psi)\) be the centered martingale
parts obtained from the exact balance and the compensated actual-contact
ledger.

\begin{lemma}[Aldous--Mitoma bounds]
\label{lem:r5-b3-tight}
For stopping times \(\tau\le T\) and \(0<h\le1\),
\[
 \mathbb E|M_{\tau+h}^\varepsilon(\phi)-M_\tau^\varepsilon(\phi)|^2
 \le C_\phi h,
\]
\[
 \mathbb E|M_{\tau+h}^\varepsilon(\phi)-M_\tau^\varepsilon(\phi)|^4
 \le C_\phi\left(h^2+\frac h{\mu_\varepsilon}\right),
\]
and the same estimates hold for \(N^\varepsilon(\psi)\).  The largest jump is
\(O(\mu_\varepsilon^{-1/2})\).  Hence every countable weighted test family is
tight, and the fields are tight in the corresponding nuclear distribution
path space.
\end{lemma}

\begin{proof}
On an interval of length \(h\), the predictable contact intensity is bounded
by the B2 constant-source moment estimate.  The quadratic variation is
therefore \(O(h)\).  Burkholder's inequality gives the fourth moment
\(O(h^2)\) plus the sum of fourth powers of jumps.  Each jump is
\(O(\mu_\varepsilon^{-1/2})\) and the expected number of jumps is
\(O(\mu_\varepsilon h)\), producing the necessary term
\(h/\mu_\varepsilon\).  Aldous' criterion gives scalar tightness; Mitoma's
theorem and the uniform weighted moments give field tightness.
\end{proof}

\begin{theorem}[Prepared joint fluctuating Boltzmann limit]
\label{thm:r5-b3-gaussian}
The pair \((\zeta^\varepsilon,\Xi^\varepsilon)\) converges to the unique
centered Gaussian solution of the linearized biased balance.  Its martingale
brackets are
\[
 d\langle M(\phi),M(\chi)\rangle_t
 =\int q\,\Delta\phi\,\Delta\chi\,dA_{f_t},
\]
\[
 d\langle M(\phi),N(\psi)\rangle_t
 =\int q\,\Delta\phi\,\psi\,dA_{f_t},
\]
\[
 d\langle N(\psi),N(\chi)\rangle_t
 =\int q\,\psi\chi\,dA_{f_t}.
\]
The drift is the Fr\'echet linearization of \(A_fq\) on the declared weighted
spaces.  Its initial covariance is first conditioned by the B1 Schur
complement and then projected to the closed balanced quotient:
\[
 \mathcal C_{\rm mc}^{\rm bal}=
 \Pi_{\rm bal}
 \left[
  \mathcal C_{\rm gc}-
  \mathcal C_{\rm gc}C^*(C\mathcal C_{\rm gc}C^*)^{-1}
  C\mathcal C_{\rm gc}
 \right]
 \Pi_{\rm bal}^*.
\]
For every finite test family its covariance equals the Hessian of the limiting
prepared pressure.
\end{theorem}

\begin{proof}
B2 analytic convergence on a larger complex chart gives convergence of all
finite-dimensional cumulants.  After central-limit scaling, cumulants of
order at least three are
\(O(\mu_\varepsilon^{1-k/2})\) and vanish.  The second derivative is exactly
the speed-normalized covariance.  Lemma~\ref{lem:r5-b3-tight} upgrades
finite-dimensional convergence to process tightness.

Differentiate the B2 marked Hamilton--Jacobi equation.  The density derivative
gives the linearized biased Boltzmann drift; differentiation in the actual
contact source gives the three displayed brackets and the linear response of
\(A_fq\).  The coefficients are continuous on the weighted chart, so the
linear Gaussian martingale problem is unique.  B1 differentiates the exact
finite-volume saddle and gives the Schur complement.  Closedness of balance
and Theorem~\ref{thm:r5-b3-kernel} give the final quotient projection.
\end{proof}

\subsection{Local information projection}

Let \(C_0\) be a finite family of continuous quotient observables.  On a
regular source chart the Hessian is positive by
Theorem~\ref{thm:r5-b3-kernel}; the mean map is therefore locally invertible.

\begin{theorem}[Regular information projection]
\label{thm:r5-b3-projection}
Every target in the mean image of a regular chart has a unique information
projection and a unique cotangent gauge class.  The minimizer and multipliers
depend analytically on the target.  No assertion is made beyond the union of
regular mean charts without a separate global steepness argument.
\end{theorem}

\begin{proof}
Apply the analytic inverse-function theorem to the pressure gradient on the
finite quotient family.  Fenchel equality and strict Hessian positivity give
the unique minimizer.  Theorem~\ref{thm:r5-b3-gauge} identifies its unique
cotangent class.
\end{proof}

\begin{theorem}[Round-five B3 closure]
\label{thm:r5-b3-main}
The Hamilton--Boltzmann cotangent is a well-typed weighted gauge quotient with
countably additive Radon dual, its covariance kernel is exactly the closed
balance-adjoint plus conserved directions, and the prepared density--actual-
contact fluctuations converge as a tight Gaussian process with all cross
brackets and the correct finite-particle increment term.
\end{theorem}

\begin{proof}
Combine Theorems~\ref{thm:r5-b3-gauge},
\ref{thm:r5-b3-duality}, \ref{thm:r5-b3-kernel},
\ref{thm:r5-b3-gaussian}, and \ref{thm:r5-b3-projection}.
\end{proof}
